\section{从数据到福利：一条可执行的工作流}
% ============================================================
% 第六讲 6.9：六步工作流、三类识别困难、检验清单
% ============================================================

\subsection{六步工作流}

前面八节的技术内容可以压缩成一条流水线。下表把每一步的输入、输出与\emph{可能失败的地方}并列，
后者比前者更值得记住。

\begin{center}
\footnotesize
\begin{tabular}{@{}c>{\raggedright\arraybackslash}p{5.2cm}>{\raggedright\arraybackslash}p{3.0cm}>{\raggedright\arraybackslash}p{4.3cm}@{}}
\toprule
步 & \multicolumn{1}{c}{操作} & \multicolumn{1}{c}{输出} & \multicolumn{1}{c}{常见失败}\\
\midrule
1 & 清洗数据、合并重复观测、定义商品与价格 & $\{(p^t,x^t)\}$ & 商品划分随数据可得性而变\\
2 & 非参数筛查：检验 GARP、算 Afriat 指数 & 是否可理性化 & 忽略角点与动态调整\\
3 & 选择函数形式（CD/CES/LES/AIDS/QUAIDS） & 参数化系统 & 形式本身已施加可积性\\
4 & 估计参数（施加加总与齐次约束） & 参数估计值 & 价格内生、测量误差\\
5 & 检验对称性与曲率 & 检验统计量、$p$ 值 & 把拒绝直接读成"理论错"\\
6 & 由 $e$ 计算 CV/EV，做敏感性分析 & 福利度量及区间 & 忽略加总与效用可比\\
\bottomrule
\end{tabular}
\end{center}
\figexplan{六步工作流。第 2 步与第 3 步是互斥的两条路：若走非参数路线，第 3--5 步可以跳过，但第 6 步无法执行（没有 $e$）；若走参数路线，第 2 步的 GARP 检验仍值得做一次作为设定诊断。表中第四列列出的是每一步最容易被忽视的失败模式，它们大多不是估计技术问题，而是"数据生成过程与模型假设不匹配"的问题。}

\subsection{第 2 步：非参数筛查怎么做}

GARP 检验的计算量随 $T$ 增长但不是瓶颈，真正的难点在两处。

\notehead{其一，价格与数量的口径必须一致。}
观测必须是"同一个决策单位、同一时期、面对同一组价格"的选择。用年度均价配年度总消费，
等于假定年内价格不变；用地区均价配家庭消费，等于假定家庭面对的是地区均价。
两种做法都会人为地\emph{减少}价格变异，从而让 GARP 更容易通过——
\textbf{检验的功效被数据构造悄悄削弱了}。这一点在论文中应当明说。

\notehead{其二，角点与耐用品的处理。}
\S 6.2 的 (A4) 假定内点，而真实数据中大量商品是零消费（不买酒、不买耐用品）。
非参数方法本身不需要内点假设，但"零消费"会使 $R^0$ 的判定失去意义（$p^t\cdot x^s$ 中缺失的分量无法补）。
常见处理是加总成大类（食品、衣着、居住……）以降低零消费比例，代价是商品数 $n$ 变小，
可检验内容随之减少（\S 6.3 的表）。

\subsection{第 4--5 步：把拒绝读对}

\begin{remark}[三种拒绝，三种解释]
设对称性检验被拒绝。至少有三种互不排斥的解释，而数据本身无法区分它们：
\begin{enumerate}
\item \textbf{偏好异质}：样本混合了参数不同的家庭，混合后的"平均需求"不是任何个体的需求（\S 6.7）；
\item \textbf{动态因素}：习惯形成、耐用品存量调整，使当期需求依赖历史，静态模型被误设；
\item \textbf{测量误差}：价格指数与数量的构造误差，在对称性检验中不抵消而累积。
\end{enumerate}
因此规范的表述是"在某某设定下对称性被拒绝"，而不是"消费者理论被证伪"。
反过来，把约束强行施加后再报告"模型通过检验"则是循环论证——
检验的对象已经被假设掉了。可行的折中是：报告无约束估计、约束估计与两者的差异，
并讨论差异是否集中在特定商品或特定收入组。
\end{remark}

\subsection{第 6 步：福利计算的三个必答问题}

在写出任何 CV/EV 数字之前，有三个问题必须先有答案：

\begin{enumerate}
\item \textbf{对象是谁？} 是个体消费者、某类家庭的代表，还是"代表性消费者"？
      三者的福利含义不同，且只有第一种有明确的偏好基础。
\item \textbf{价格变化的路径是什么？} 单一价格变动时 CV/EV 无路径依赖（命题~\ref{prop:no-path}），
      多价格同时变动时也无——但如果用 $\Delta\mathrm{CS}$ 做近似，路径必须交代。
\item \textbf{分配效应有多大？} 若政策同时改变相对价格与收入分配，
      代表性消费者的 CV 可能与人头加总的福利变化符号相反（\S 6.7 的注）。
      至少应报告分收入组的 CV，看方向是否一致。
\end{enumerate}

\subsection{一个可复用的检验清单}

把本章的技术条件整理成一张可以直接贴进论文附录的清单：

\begin{center}
\begin{tabular}{lll}
\toprule
检验对象 & 原假设 & 拒绝的含义\\
\midrule
零次齐次 & $\sum_j\gamma_{ij}=0$ & 存在货币幻觉或价格指数误设\\
Walras 律 & $\sum_ip_ix_i=w$ & 储蓄、耐用品存量或数据口径问题\\
对称性 & $S=S^{\top}$（或 $\gamma_{ij}=\gamma_{ji}$） & 异质、动态或测量误差\\
负半定性 & $S$ 的特征值 $\le0$ & 补偿需求向上倾斜，理论被实质违反\\
GARP & 无循环严格偏好 & 数据不可由任何局部非饱和效用函数生成\\
\bottomrule
\end{tabular}
\end{center}
\figexplan{五个检验的对象与解读。前两行是会计恒等式，拒绝通常指示数据问题而非理论问题；后三行才是实质检验。注意检验的次序有讲究：应先查前两行（数据是否自洽），再查对称与曲率（理论是否成立），最后查 GARP（是否需要放弃参数化）。顺序颠倒会导致把数据错误误判为理论失败。}
